\section{Installation on the Raspberry Pi}

\subsection{System Requirements}

Automated image acquisition requires a Raspberry Pi 4 or newer, a Raspberry Pi
Camera Module v3, Raspberry Pi OS Bookworm or newer, a working system clock,
and sufficient storage. Phenara runs as separate scheduler and web-interface
systemd services.

\subsection{Installing Phenara}

Phenara does not require a dedicated Linux account or fixed installation path.
Clone the repository as the user that should run the services and invoke the
installer from that checkout.

\begin{verbatim}
git clone https://github.com/kfreinders/phenara.git
cd phenara
./deploy/install.sh
\end{verbatim}

The installer performs the following operations:

\begin{itemize}
  \item installs the required apt packages on Raspberry Pi OS or Debian;
  \item creates a project-local \texttt{.venv} with access to Picamera2;
  \item installs the declared Python dependencies;
  \item installs and builds the React frontend;
  \item creates the runtime and capture directories;
  \item writes \texttt{/etc/phenara/phenara.env};
  \item installs, enables, and starts both systemd services.
\end{itemize}

It is safe to rerun the installer after updating the repository. The
\texttt{--skip-system-packages} option is useful on non-Debian development
machines, while \texttt{--no-start} installs and enables the units without
starting them. Sample-image development mode is unavailable by default. It can
be included explicitly with \texttt{--enable-development-mode}.

\subsection{Shared Configuration}

All paths are defined by \texttt{phenara/config.py}. The installed services
load the same environment file, containing:

\begin{verbatim}
PHENARA_ROOT
PHENARA_RUNTIME_DIR
PHENARA_CAPTURE_DIR
PHENARA_DEVELOPMENT_AVAILABLE
PHENARA_VENV_DIR
PHENARA_PYTHON
PHENARA_TIMEZONE
PHENARA_GUI_HOST
PHENARA_GUI_PORT
\end{verbatim}

Without installation, defaults are derived from the current source checkout.
Runtime data, captures, and the virtual environment are stored in
\texttt{runtime/}, \texttt{captures/}, and \texttt{.venv/}, respectively.

\subsection{Checking the Services}

\begin{verbatim}
sudo systemctl status phenara-scheduler phenara-gui
sudo journalctl -u phenara-scheduler -u phenara-gui -f
\end{verbatim}

Both services start at boot and restart automatically after a failure. The web
interface listens on port 8000 by default.
